\documentclass{beamer}
\usepackage{amsmath, booktabs}
\usetheme{Madrid}


\title{The Importance of Economic Growth in the Evolution of the Standard of Living}
\subtitle{Contemporary Economic Issues: ECON 204}
\author{Muhebullah Karimzada}
\date{Winter 2025}







\begin{document}

\frame{\titlepage}

\begin{frame}{Introduction}
    \begin{itemize}
        \item \textbf{Definition of Economic Growth}: Increase in a nation’s output of goods and services, measured by GDP.
        \item \textbf{Definition of Standard of Living}: The level of wealth, comfort, material goods, and necessities available to a population.
        \item \textbf{Key Question}: How has economic growth shaped human well-being over time?
        \item \textbf{Explanation}: Economic growth enables societies to produce more goods and services, improving overall quality of life. Without sustained growth, economies stagnate, leading to high unemployment and poor living conditions.
    \end{itemize}
\end{frame}

\begin{frame}{The Link Between Economic Growth and Living Standards}
    \begin{itemize}
        \item Higher incomes lead to better access to goods and services.
        \item Economic expansion enables greater investment in healthcare, education, and infrastructure.
        \item Growth fosters innovation, enhancing productivity and efficiency.
        \item \textbf{Example:} Post-WWII U.S. economic boom led to higher wages, homeownership, and increased consumer goods access.
        \item \textbf{Example (Canada):} Canada's resource-driven growth has resulted in high living standards, universal healthcare, and social programs.
    \end{itemize}
\end{frame}

\begin{frame}{Historical Perspective on Economic Growth and Living Standards}
    \begin{itemize}
        \item \textbf{Pre-Industrial Revolution}: Slow, stagnant growth with widespread poverty.
        \item \textbf{Industrial Revolution (1750–1900)}: Technological advancements raised productivity and incomes.
        \item \textbf{20th Century}: Global GDP per capita surged, reducing poverty and improving health and education.
        \item \textbf{21st Century}: Digital and AI-driven economies continue shaping modern living standards.
        \item \textbf{Example (Canada)}: Canadian Pacific Railway expansion in the late 19th century was crucial for national economic growth and industrialization.
    \end{itemize}
\end{frame}

\begin{frame}{Economic Growth and Health Outcomes}
    \begin{itemize}
        \item Higher GDP allows investment in healthcare infrastructure.
        \item Advances in medicine and public health increase life expectancy.
        \item \textbf{Example}: The Green Revolution increased agricultural productivity, reducing malnutrition.
        \item \textbf{Example (Canada)}: Universal healthcare, funded by economic growth, contributes to high life expectancy.
    \end{itemize}
\end{frame}

\begin{frame}{Economic Growth and Education}
    \begin{itemize}
        \item Increased national income supports universal education access.
        \item Higher literacy rates improve labor productivity and innovation.
        \item \textbf{Example}: South Korea’s education-driven industrialization led to advanced economic status.
        \item \textbf{Example (Canada)}: High post-secondary education rates due to government-funded institutions.
    \end{itemize}
\end{frame}

\begin{frame}{Infrastructure and Economic Growth}
    \begin{itemize}
        \item Growth enables investment in transportation, energy, and communication.
        \item Modern infrastructure increases efficiency and connectivity.
        \item \textbf{Example (Canada)}: Trans-Canada Highway and urban transit projects have supported economic growth.
    \end{itemize}
\end{frame}

\begin{frame}{The Role of Technological Innovation}
    \begin{itemize}
        \item Innovation is both a driver and consequence of economic growth.
        \item Technological advances improve productivity and reduce costs.
        \item \textbf{Example}: E-commerce platforms like Amazon revolutionized global commerce.
        \item \textbf{Example (Canada)}: Canada’s AI and clean energy sectors drive economic growth, especially in tech hubs like Toronto and Vancouver.
    \end{itemize}
\end{frame}

\begin{frame}{Income Inequality and Economic Growth}
    \begin{itemize}
        \item While growth improves living standards, benefits are not always evenly distributed.
        \item Policies such as progressive taxation and social safety nets mitigate inequality.
        \item \textbf{Example (Canada)}: Canada’s progressive taxation and social programs (e.g., Child Benefit, EI) support equitable distribution.
    \end{itemize}
\end{frame}

\begin{frame}{Challenges to Sustaining Economic Growth}
    \begin{itemize}
        \item Environmental constraints (climate change, resource depletion).
        \item Rising debt and financial instability.
        \item Job displacement from automation and AI.
        \item \textbf{Example (Canada)}: Balancing economic growth with environmental sustainability, particularly in the energy sector.
    \end{itemize}
\end{frame}

\begin{frame}{Future Prospects for Economic Growth and Living Standards}
    \begin{itemize}
        \item Green and sustainable growth as a new paradigm.
        \item The role of AI and digital transformation in productivity.
        \item Global cooperation in addressing inequality and climate change.
        \item \textbf{Example (Canada)}: Investments in clean energy (e.g., hydroelectric power, carbon capture) to ensure sustainable growth.
    \end{itemize}
\end{frame}

\begin{frame}{Conclusion}
    \begin{itemize}
        \item Economic growth has been the primary driver of rising living standards.
        \item Sustainable and inclusive growth is crucial for long-term prosperity.
        \item Policymakers must balance economic expansion with social and environmental concerns.
    \end{itemize}
\end{frame}

\begin{frame}{Discussion Questions}
    \begin{enumerate}
        \item Can economic growth occur without environmental degradation?
        \item How can governments ensure the benefits of growth are equally distributed?
        \item What are the biggest threats to future economic growth?
    \end{enumerate}
\end{frame}

\begin{frame}{Exercise 1: Doubling Time Formula (Question)}
    \begin{itemize}
        \item Use the Rule of 70 to estimate the time required for GDP to double given different growth rates:
        \begin{enumerate}
            \item If the growth rate is 2\%, how many years will it take for GDP to double?
            \item If the growth rate is 3.5\%, how many years will it take?
            \item What growth rate is needed for GDP to double in 25 years?
        \end{enumerate}
    \end{itemize}
\end{frame}

\begin{frame}{Exercise 1: Solution}
    \begin{itemize}
        \item The Rule of 70 formula:
        \begin{equation*}
            \text{Doubling Time} \approx \frac{70}{\text{Growth Rate (\%)}}
        \end{equation*}
        \item Given a growth rate of 2\%:
        \begin{equation*}
            \text{Doubling Time} = \frac{70}{2} = 35 \text{ years}
        \end{equation*}
        \item Given a growth rate of 3.5\%:
        \begin{equation*}
            \text{Doubling Time} = \frac{70}{3.5} \approx 20 \text{ years}
        \end{equation*}
        \item To double in 25 years:
        \begin{equation*}
            \text{Growth Rate} = \frac{70}{25} \approx 2.8\%
        \end{equation*}
    \end{itemize}
\end{frame}

\begin{frame}{Exercise 2: Future GDP Per Capita (Question)}
    \begin{itemize}
        \item Using the GDP growth formula, calculate the future GDP per capita for two countries:
        \begin{enumerate}
            \item Country A starts at $40,000$ with a growth rate of 1.5\%. What is the GDP per capita after 20 years?
            \item Country B starts at $15,000$ with a growth rate of 5\%. What is the GDP per capita after 20 years?
        \end{enumerate}
    \end{itemize}
\end{frame}

\begin{frame}{Exercise 2: Solution}
    \begin{itemize}
        \item The growth formula:
        \begin{equation*}
            GDP_{t} = GDP_{0} \times (1 + g)^t
        \end{equation*}
        \item For Country A ($40,000$, 1.5\% growth, $t=20$):
        \begin{equation*}
            GDP_{20} = 40,000 \times (1.015)^{20} \approx 54,000
        \end{equation*}
        \item For Country B ($15,000$, 5\% growth, $t=20$):
        \begin{equation*}
            GDP_{20} = 15,000 \times (1.05)^{20} \approx 39,800
        \end{equation*}
    \end{itemize}
\end{frame}

\begin{frame}{Key Takeaways}
    \begin{itemize}
        \item Small differences in growth rates have large effects over time.
        \item Faster growth leads to a higher standard of living.
        \item Policies that enhance productivity and growth are crucial for prosperity.
    \end{itemize}
\end{frame}

\end{document}

